\subsection{应用：I. 连续函数到测度空间的延拓}

设 T 为局部紧空间， F 为拟完备的 Hausdorff 局部凸空间， f 为 T 到 F 中的连续映射；若 \(\mu\) 是 T 上的正测度，具有紧支集 S ，则 \(\mathbf{f}(S)\) 是紧的；于是 \(\mathbf{f}(S)\) 的闭凸包络是紧的（TVS, III, §1, No. 6），因此 f 是标量 \(\mu\)-可积的且 \(\int f d\mu \in F\) （No. 2, 命题 5 的推论）。现在若 \(\lambda\) 是任一具有紧支集的实测度，则 \(\lambda^{+}\) 与 \(\lambda^{-}\) 都是具有紧支集的正测度；令 \(\int f d\lambda = \int f d\lambda^{+} - \int f d\lambda^{-}\) ，即可立即验证（利用关系式 \((\lambda + \mu)^{+} + \lambda^{-} + \mu^{-} = \lambda^{+} + \mu^{+} + (\lambda + \mu)^{-}\) ） \(\lambda \mapsto \int f d\lambda\) 是由 T 上具有紧支集的测度所成的空间 \(\mathcal{C}'(T)\) 到局部凸空间 F 中的线性映射。

现在我们注意到，空间 \(\mathcal{C}^{\prime}(\mathrm{T})\) 可以等同于 T 上连续数值函数所成的空间 \(\mathcal{C}(\mathrm{T})\) 的对偶（其记法即由此而来），但须 \(\mathcal{C}(\mathrm{T})\) 赋以紧收敛拓扑（在本目及下一目中我们总作此假设）：因为一方面已知（Ch. IV, §4, No. 8, 命题 14） T 上可以延拓为 \(\mathcal{C}(\mathrm{T})\) 上的连续线性形式的测度恰是具有紧支集的测度；反过来，连续线性形式在 \(\mathcal{K}(T)\) 上的限制（该形式定义在 \(\mathcal{C}(T)\) 上）是一个测度（因为 \(\mathcal{K}(T)\) 的拓扑细于 \(\mathcal{C}(T)\) 的拓扑所诱导的拓扑）。

命题 14. --- 设 T 为局部紧空间， F 为拟完备的 Hausdorff 局部凸空间， f 为 T 到 F 中的连续映射。若 T 上具有紧支集的测度所成的空间 \(\mathcal{C}'(T)\) 赋以在 \(\mathcal{C}(T)\) 的紧子集上一致收敛的拓扑，则映射 \(\lambda \mapsto \int \mathbf{f} d\lambda\) 是唯一的连续线性映射 \(\widetilde{\mathbf{f}}\) （ \(\mathcal{C}'(T)\) 到 F 中），使得 \(\widetilde{\mathbf{f}}(\varepsilon_t) = \mathbf{f}(t)\) 对每个 \(t \in T\) 成立。

为确立延拓的唯一性，只需看出点测度 \(\varepsilon_{t}\) 在 \(\mathcal{C}'(T)\) 中构成一个完全集；由于 \(\mathcal{C}'(T)\) 的对偶是 \(\mathcal{C}(T)\) （TVS, IV, §1, No. 1, 定理 1），只需注意到每个函数 \(g \in \mathcal{C}(T)\) 若与所有测度 \(\varepsilon_{t}\) 正交，则按定义为 0（TVS, IV, §1, No. 2, 命题 2）。

现在证明 \(\lambda \mapsto \int \mathbf{f}d\lambda\) 连续。令 \(V\) 为 \(F\) 中 0 的闭、平衡、凸邻域；只需证明存在相对紧子集 \(L\) （ \(\mathcal{C}(T)\) 中的），使得关系 \(\lambda \in L^{\circ}\) 与 \(\mathbf{z}' \in V^{\circ}\) 蕴含 \(\left|\left\langle \int \mathbf{f}d\lambda, \mathbf{z}' \right\rangle \right| \leqslant 1\) ，或等价地 \(\left|\int \langle \mathbf{f}, \mathbf{z}' \rangle d\lambda \right| \leqslant 1\) 。为此，我们将证明当 \(\mathbf{z}'\) 跑遍 \(V^{\circ}\) 时，集 \(L\) ——由数值函数 \(\langle \mathbf{f}, \mathbf{z}' \rangle\) 组成——在 \(\mathcal{C}(T)\) 中相对紧。由于 \(V^{\circ}\) 对 \(\sigma(F', F)\) 有界，数 \(|\langle \mathbf{f}(t), \mathbf{z}' \rangle|\) （对固定的 \(t \in T\) 与 \(\mathbf{z}'\) 跑遍 \(V^{\circ}\) ）的上确界是有限的；根据 Ascoli 定理（GT, X, §2, No. 5, 定理 2 的推论 2），因此只需证明诸 \(\langle \mathbf{f}, \mathbf{z}' \rangle\) （ \(\mathbf{z}' \in V^{\circ}\) ）所成之集等度连续。但是，对每个 \(t_0 \in T\) 与每个 \(\delta > 0\) ，按假设存在一个邻域 \(W\) （ \(t_0\) 在 \(T\) 中的邻域），使得 \(\mathbf{f}(t) - \mathbf{f}(t_0) \in \delta V\) （对所有 \(t \in W\) ）；由此推出 \(|\langle \mathbf{f}(t), \mathbf{z}' \rangle - \langle \mathbf{f}(t_0), \mathbf{z}' \rangle| \leqslant \delta\) （对所有 \(t \in W\) 与所有 \(\mathbf{z}' \in V^{\circ}\) ），证毕。

注记. --- 1) 映射 \(t \mapsto \varepsilon_t\) 是 T 到空间 \(\mathcal{C}'(\mathrm{T})\) 中的同胚；因为，若 L 是 \(\mathcal{C}(\mathrm{T})\) 的紧子集，且 \(t_0 \in \mathrm{T}\) ，则存在（GT, X, §2, No. 5, 定理 2 的推论 3） \(t_0\) 的一个邻域 W ，使得 \(|g(t) - g(t_0)| \leqslant 1\) （对每个 \(t \in \mathrm{W}\) 与每个函数 \(g \in \mathrm{L}\) ），因此 \(\varepsilon_t - \varepsilon_{t_0} \in \mathrm{L}^\circ\) （对 \(t \in \mathrm{W}\) ），这就证明了 \(t \mapsto \varepsilon_t\) 的连续性（参见 Ch. IV, §4, No. 8, 命题 15）；此外已知逆映射对淡拓扑已经连续（Ch. III, §1, No. 9, 命题 13），从而对在 \(\mathcal{C}(\mathrm{T})\) 的紧子集上一致收敛的拓扑更是如此。如果进而把 T 与它在 \(\mathcal{C}'(\mathrm{T})\) 中的像（借助 \(t \mapsto \varepsilon_t\) ）等同起来，就可以说 \(\lambda \mapsto \int f d\lambda\) 是 f 到线性映射的唯一连续延拓。

\begin{enumerate}
\def\labelenumi{\arabic{enumi})}
\setcounter{enumi}{1}
\tightlist
\item
  注意，在 \(\lambda\mapsto\int f d\lambda\) 的连续性的证明中，我们没有用到 F 拟完备这一事实。因此命题 14 的结论在没有这一假设时仍然成立，只要另外已知 \(\int f d\mu\in F\) 对每个具有紧支集的正测度 \(\mu\) 成立。
\end{enumerate}

现在设 \(\mathbf{f}(\mathrm{T})\) 是 \(\mathrm{F}\) 的有界子集。于是，对每个有界正测度 \(\mu\) （ \(\mathrm{T}\) 上的）， \(\mathbf{f}\) 是标量 \(\mu\)-可积的，且 \(\int \mathbf{f}d\mu \in \mathrm{F}\) （No. 2, 命题 8）。若 \(\lambda\) 是 T 上任一有界实测度，则 \(\lambda^{+}\) 与 \(\lambda^{-}\) 都有界，并且立即可以看出，如上定义的 \(\lambda \mapsto \int \mathbf{f} d\lambda\) 是 T 上有界测度所成的空间 \(\mathcal{M}^1(\mathrm{T})\) 到局部凸空间 F 中的线性映射，它显然延拓了映射 \(\lambda \mapsto \int \mathbf{f} d\lambda\) （ \(\mathcal{C}'(\mathrm{T})\) 到 F 中）。

命题 15. --- 设 T 为局部紧空间， F 为拟完备的 Hausdorff 局部凸空间， f 为 T 到 F 中的连续映射，使得 \(\mathbf{f}(\mathrm{T})\) 有界。若 \(\mathcal{M}^{1}(\mathrm{T})\) 赋以它的 Banach 空间拓扑，则线性映射 \(\lambda \mapsto \int f d\lambda\) （ \(\mathcal{M}^{1}(\mathrm{T})\) 到 F 中）是连续的。

对 F 中 0 的每个闭、平衡、凸邻域 V ，存在 \(\rho > 0\) 使 \(\mathbf{f}(\mathrm{T}) \subset \rho\mathrm{V}\) ；于是 \(\mathbf{f}(\mathrm{T})\) 的闭、平衡、凸包络 B 含于 \(\rho V\) ，且由假设它是完备的。这时若 \(\|\lambda\| \leqslant 1/\rho\) ，则由第 2 目命题 8 及关系式 \(\|\lambda\| = \lambda^{+}(\mathrm{T}) + \lambda^{-}(\mathrm{T})\) 得 \(\int f d\lambda \in B/\rho \subset V\) 。

